% !TeX root = ../main.tex
%% ===========================================================================
%%  Using This Template
%%  This chapter explains how to use the template and gives advice on how to
%%  write a Master's thesis. It is only included when compiling
%%  template-guide.tex (which produces template-guide.pdf); main.tex, your
%%  thesis, leaves it out. You can delete this file if you do not need it.
%% ===========================================================================
\chapter{Using This Template}\label{ch:guide}

This chapter explains how to use this \LaTeX{} template to write a Master's thesis
(\textit{\gls{tfm}}) in English at a Spanish university. The template is designed for
\texttt{pdflatex}, the most widely used \LaTeX{} compiler, together with \texttt{biber} for
the bibliography, and it also works with \texttt{lualatex}.

You can compile the sources with a local \LaTeX{} distribution, such as TeX Live
(\url{https://tug.org/texlive/}) or MiKTeX (\url{https://miktex.org/}), or online with
Overleaf (\url{https://www.overleaf.com}), which needs no installation. The template is an
English, Master's-level and university-independent adaptation of the TFG template created
by Francisco Charte Ojeda for the Escuela Politécnica Superior de Jaén
(\url{https://github.com/fcharte/TFGLatexTemplate}).

\begin{guidance}
Boxes like this one contain guidance about what to write in each part of the thesis. They
are only visible while you write: add the \texttt{final} option to the
\texttt{\textbackslash documentclass} line of \texttt{main.tex} and all of them
disappear from the PDF, together with the to-do notes. This chapter itself is not part of
your thesis: it only appears when compiling \texttt{template-guide.tex}, the document you
are reading, and never in \texttt{main.tex}.
\end{guidance}

\section{Before you start: your university's regulations}\label{sec:regulations}

In Spain, Master's degrees are regulated by Royal Decree 822/2021~\cite{rd822-2021}, which
establishes that every Master's degree ends with a thesis worth between 6 and 30
\gls{ects} credits that must be defended in public. Beyond that common framework, each
university, and often each school or Master's programme, sets its own rules: length,
structure, formatting, the language in which the thesis can be written, whether an abstract
in Spanish or in a co-official language is required, the declarations that must be signed,
and how the thesis is submitted and archived.

This template implements the requirements that are most common across Spanish public
universities and lets you adapt the rest. Before writing, read the regulations of your
university and your Master's programme, and check in particular:

\begin{itemize}
  \item whether you need \emph{permission} to write the thesis in English, and whether the
        title, abstract, keywords or conclusions must also be written in Spanish or in a
        co-official language (option \texttt{secondlanguage});
  \item whether there is an \emph{official cover} that must be used (field
        \texttt{\textbackslash OfficialCover}) and which logos may appear;
  \item the \emph{formatting} rules: margins, typeface, size and line spacing;
  \item the \emph{declarations} required: originality, supervisor's approval, use of
        generative artificial intelligence, alignment with the \glspl{sdg};
  \item the \emph{maximum length} and the citation style expected in your field.
\end{itemize}

\section{Getting started}\label{sec:getting-started}

\subsection{Files of the template}\label{sec:files}

\Cref{tab:files} describes the files and folders of the template. In most cases you only
need to edit \texttt{config.tex}, the files in \texttt{frontmatter/} and \texttt{chapters/},
and \texttt{references.bib}.

\begin{table}[htbp]
  \centering
  \small
  \caption{Files and folders of the template.}
  \label{tab:files}
  \begin{tabularx}{\textwidth}{@{}>{\ttfamily}l L@{}}
    \toprule
    \textnormal{\textbf{File or folder}} & \textbf{Purpose} \\
    \midrule
    main.tex         & Main document of your thesis: class options and order of its parts. \\
    template-guide.tex & This user guide: your thesis plus this chapter. \\
    config.tex       & Your data: title, author, supervisors, university, logos, licence,
                       optional pages. \\
    tfmthesis.cls    & The class that defines the layout. You do not need to edit it. \\
    frontmatter/     & Dedication, acknowledgements, abstracts, AI statement and acronyms. \\
    chapters/        & One file per chapter. \\
    appendices/      & One file per appendix. \\
    figures/         & Images and plots (\texttt{figures/src/} holds the scripts that
                       generate them, \texttt{figures/logos/} the logos). \\
    tables/          & Long tables kept in separate files. \\
    references.bib   & Bibliographic references in BibLaTeX format. \\
    presets/         & Settings that reproduce the style of specific universities. \\
    latexmkrc        & Compilation settings for \texttt{latexmk} and Overleaf. \\
    \bottomrule
  \end{tabularx}
\end{table}

\subsection{Setting your data}\label{sec:config}

Open \texttt{config.tex} and fill in your data, following the comments in the file.
\Cref{lst:config} shows its most important fields. These data are used automatically on the
cover, on the declarations, in the page headers and in the metadata of the PDF file.

\begin{lstlisting}[language={[LaTeX]TeX}, caption={Main fields of \texttt{config.tex}.}, label=lst:config]
\ThesisTitle{Title of Your Master's Thesis}
\ThesisShortTitle{Short title for the page headers}
\ThesisTitleSecondLanguage{Título del trabajo en español}
\ThesisAuthor{First-name Surname1 Surname2}
\Supervisor{Dr. First-name Surname}{Department of ..., University of ...}
\Supervisor[Co-supervisor]{Dr. First-name Surname}{Company or institute}
\University{Universidad de ...}
\Faculty{School of ...}
\Degree{Master's Degree in ...}
\AcademicYear{2025/2026}
\SubmissionDate{June 2026}
\UniversityLogo{figures/logos/university-logo.pdf}
\definecolor{tfmaccent}{HTML}{1F4E79}
\end{lstlisting}

The \texttt{\textbackslash Supervisor} command can be repeated as many times as needed, and
its optional argument sets the role shown on the cover (Supervisor, Co-supervisor, External
supervisor\dots). If a field is left empty, it disappears from the document. In your own
text you can reuse these data with \texttt{\textbackslash thesistitle},
\texttt{\textbackslash thesisauthor}, \texttt{\textbackslash thesisdegree},
\texttt{\textbackslash thesisuniversity} and \texttt{\textbackslash thesissupervisors}.

\subsection{Class options}\label{sec:options}

The appearance of the document is controlled with the options of the
\texttt{\textbackslash documentclass} command in \texttt{main.tex}, summarised in
\cref{tab:options}. For example, the final version of a thesis written in a sans-serif
typeface, to be printed, with the abstract also in Catalan and references in APA style,
would start with:

\begin{lstlisting}[language={[LaTeX]TeX}, caption={Example of class options.}, label=lst:options]
\documentclass[print, final, font=sans, secondlanguage=catalan,
               bibstyle=apa]{tfmthesis}
\end{lstlisting}

\begin{table}[htbp]
  \centering
  \small
  \caption{Options of the \texttt{tfmthesis} class. Default values are shown in bold.}
  \label{tab:options}
  \begin{tabularx}{\textwidth}{@{}>{\ttfamily}l L@{}}
    \toprule
    \textnormal{\textbf{Option}} & \textbf{Effect} \\
    \midrule
    print          & Two-sided layout for a printed and bound copy: wider inner margin,
                     chapters start on odd pages and links are black. Without it the
                     document is optimised for the screen (\textbf{digital}). \\
    final          & Hides the guidance boxes and the to-do notes. \\
    review         & Adds line numbers and a \enquote{DRAFT} watermark, so that your
                     supervisors can refer to specific lines in their comments. \\
    font=          & \textbf{\texttt{serif}} (Times-like), \texttt{sans} (Helvetica/Arial-like)
                     or \texttt{modern} (Latin Modern, the classic \LaTeX{} look). \\
    secondlanguage=& Language of the second abstract: \textbf{\texttt{spanish}},
                     \texttt{catalan}, \texttt{valencian}, \texttt{galician},
                     \texttt{basque} or \texttt{none}. \\
    bibstyle=      & Any BibLaTeX style: \textbf{\texttt{numeric-comp}}, \texttt{ieee},
                     \texttt{apa}, \texttt{authoryear}, \texttt{nature}\dots \\
    english=       & \textbf{\texttt{british}} or \texttt{american} spelling and hyphenation. \\
    uppercase      & Chapter titles in capital letters. \\
    11pt           & Smaller body text (\textbf{12pt} by default). \\
    \bottomrule
  \end{tabularx}
\end{table}

\subsection{Optional pages}\label{sec:optional-pages}

\Cref{tab:optional-pages} lists the optional pages of the front matter and how to include or
remove them. The order in which they appear is set in \texttt{main.tex}.

\begin{table}[htbp]
  \centering
  \small
  \caption{Optional pages of the front matter.}
  \label{tab:optional-pages}
  \begin{tabularx}{\textwidth}{@{}l L@{}}
    \toprule
    \textbf{Page} & \textbf{How to control it} \\
    \midrule
    Official cover          & \texttt{\textbackslash OfficialCover\{file.pdf\}} in \texttt{config.tex} \\
    Licence / copyright     & Toggle \texttt{tfmLicencePage}; licence chosen with
                              \texttt{\textbackslash ThesisLicence} \\
    Supervisor's approval   & Toggle \texttt{tfmApprovalPage} \\
    Declaration of originality & Toggle \texttt{tfmOriginalityPage} \\
    Dedication, AI statement, abstracts & Comment out their \texttt{\textbackslash input}
                              line in \texttt{main.tex} \\
    Lists of figures, tables, algorithms, listings and acronyms & Comment out the
                              corresponding line in \texttt{main.tex} \\
    \bottomrule
  \end{tabularx}
\end{table}

\subsection{Compiling the document}\label{sec:compiling}

The bibliography is processed with \texttt{biber}, so a full compilation needs several
steps. The simplest way is to let \texttt{latexmk} run them for you:

\begin{lstlisting}[language=bash, caption={Compiling the thesis from the command line.}, label=lst:compile]
latexmk main.tex          # builds main.pdf (pdfLaTeX + Biber)
latexmk -pvc main.tex     # recompiles automatically whenever a file changes
latexmk -c                # removes the auxiliary files
\end{lstlisting}

Overleaf, TeXstudio, VS Code with LaTeX Workshop and most editors detect all of this
automatically. In Overleaf, keep the compiler set to pdfLaTeX (Menu~$\rightarrow$ Compiler).

\paragraph{Compilation time on Overleaf} The free plan of Overleaf stops any compilation
that takes longer than a few seconds, and a thesis grows quickly. The template is tuned to
compile fast: the extra compilation passes that \texttt{latexmk} would otherwise run are
avoided (file \texttt{latexmkrc}), and microtypography, which noticeably slows down
every compilation, is only enabled with the \texttt{final} option. If your thesis starts to
time out, the following tips help:
\begin{itemize}
  \item Leave out the chapters you are not working on: comment out their
        \verb|\input| lines in \texttt{main.tex}. Cross-references to them will show
        \enquote{??} until you include them again.
  \item Use \enquote{Fast [draft]} compile mode in Overleaf (Menu~$\rightarrow$ Compile
        mode), which skips images.
  \item Keep figures as PDF or reasonably sized PNG/JPEG files: a 20~MB photograph
        slows down every compilation.
  \item Compile locally with \texttt{latexmk}, where there is no time limit, or ask your
        university whether it provides Overleaf Professional accounts.
\end{itemize}

\section{Structuring the thesis}\label{sec:structuring}

A thesis is similar to a book, so its content is organised in several levels with a
hierarchical numbering, which you can see in the table of contents. The usual levels are:

\begin{enumerate}[label={}, leftmargin=1em]
  \item 1. Chapter (\verb|\chapter{Title}|)
  \begin{enumerate}[label={}, leftmargin=1.5em]
    \item 1.1. Section (\verb|\section{Title}|)
    \begin{enumerate}[label={}, leftmargin=1.5em]
      \item 1.1.1. Subsection (\verb|\subsection{Title}|)
      \begin{enumerate}[label={}, leftmargin=1.5em]
        \item 1.1.1.1. Subsubsection (\verb|\subsubsection{Title}|)
      \end{enumerate}
    \end{enumerate}
  \end{enumerate}
\end{enumerate}

More than four levels are not recommended. Subsubsections are numbered but, to keep the table
of contents short, they are not listed in it; add \verb|\setcounter{tocdepth}{3}| to
\texttt{config.tex} to include them.

\subsection{Chapters}\label{sec:chapters}

Each chapter is written in a separate \texttt{.tex} file inside \texttt{chapters/}, which is
included in \texttt{main.tex} with \verb|\input{chapters/file}|. The first line of each
chapter defines its title and a label to refer to it:

\begin{lstlisting}[language={[LaTeX]TeX}, caption={Beginning of a chapter.}, label=lst:chapter]
\chapter{Materials and Methods}\label{ch:methods}
\end{lstlisting}

Write the titles in title case; the numbering is automatic. If your university requires
chapter titles in capitals, use the \texttt{uppercase} class option instead of typing them
in capitals, so that the table of contents and the headers stay consistent.

\subsection{Sections and subsections}\label{sec:sections}

Plan the sections of a chapter before writing it: think about what you want to tell and how
to split it into cohesive parts, each of which becomes a section. Long sections are divided
into subsections. In general, several pages of text without any division make the document
hard to follow.

\subsection{Paragraphs}\label{sec:paragraphs}

Paragraphs are separated by a blank line in the source. The text of a paragraph is written
continuously, without pressing Enter: \LaTeX{} breaks lines and hyphenates words
automatically. Do not use \verb|\\| or \verb|\newline| to separate paragraphs. It is good
practice to have at least one paragraph between two consecutive headings, so that a section
never starts immediately after the title of its chapter.

\subsection{Cross-references}\label{sec:cross-references}

Referring to other parts of the document (chapters, sections, figures, tables, equations\dots)
takes two steps:

\begin{enumerate}
  \item Put a label right after the element, with \verb|\label{key}|. Using a prefix that
        identifies the type of element (\texttt{ch:}, \texttt{sec:}, \texttt{fig:},
        \texttt{tab:}, \texttt{eq:}, \texttt{lst:}, \texttt{alg:}) makes keys easier to
        remember.
  \item Refer to it with \verb|\cref{key}|, which inserts both the type of element and its
        number, e.g.\ \enquote{\cref{sec:chapters}}. Use \verb|\Cref| at the beginning of a
        sentence.
\end{enumerate}

The \verb|\nameref| command inserts the title of the element and \verb|\cpageref| its page.
For example\footnote{The \texttt{\textasciitilde} character is a non-breaking space: it
keeps the words on both sides on the same line.}:

\begin{lstlisting}[language={[LaTeX]TeX}, caption={Cross-references.}, label=lst:refs]
\Cref{sec:chapters}, \nameref{sec:chapters}, on \cpageref{sec:chapters},
explains how to split the thesis into chapters.
\end{lstlisting}

produces: \enquote{\Cref{sec:chapters}, \nameref{sec:chapters}, on \cpageref{sec:chapters},
explains how to split the thesis into chapters.} All references are hyperlinks in the PDF and
are updated automatically when the document changes.

\section{Text formatting}\label{sec:formatting}

The styles of headings and paragraphs are predefined. In the text you will mainly use three
variants:

\begin{itemize}
  \item \textbf{Bold}, with \verb|\textbf{...}|: use it sparingly to \textbf{highlight
        something important}; if there is too much bold text, it loses its effect.
  \item \textit{Italics}, with \verb|\textit{...}| or \verb|\emph{...}|: for terms in other
        languages, e.g.\ \textit{Trabajo Fin de Máster}, for titles of works and
        occasionally for emphasis.
  \item \texttt{Monospace}, with \verb|\texttt{...}|: for names of variables, functions,
        files and anything related to code, e.g.\ the \texttt{print()} function of Python.
\end{itemize}

\LaTeX{} allows colours, font sizes and many other effects, but a thesis is a formal
document: use them only when they convey information. Quotation marks are written with
\verb|\enquote{...}|, which produces the correct \enquote{curly quotes} for the language.

\subsection{Footnotes}\label{sec:footnotes}

Footnotes\footnote{Footnotes are numbered automatically within each chapter.} are useful to
clarify something or add complementary information without breaking the flow of the text.
They are inserted with \verb|\footnote{Text}| right where the number must appear. Like bold
text, do not abuse them.

\subsection{Lists}\label{sec:lists}

Lists make enumerations and procedures with several steps easier to read. Use a numbered
list (\texttt{enumerate}) when the order matters, and a bulleted list (\texttt{itemize})
otherwise. In \LaTeX{} lists are \emph{environments}: blocks of content that start with
\verb|\begin{name}| and end with \verb|\end{name}|. Each element starts with \verb|\item|, as
shown in \cref{lst:itemize}. Lists can be nested, and the indentation, bullets and
numbering adjust automatically.

\begin{lstlisting}[language={[LaTeX]TeX}, caption={A bulleted list with two items.}, label=lst:itemize]
\begin{itemize}
  \item \texttt{itemize}: produces a bulleted list
  \item \texttt{enumerate}: produces a numbered list
\end{itemize}
\end{lstlisting}

\section{Figures}\label{sec:figures}

A thesis contains many kinds of figures: diagrams of a system, plots of results,
screenshots, photographs\dots{} Keep the following rules in mind:

\begin{itemize}
  \item Every figure has a number and a \emph{caption} that describes clearly what it shows.
        A long caption is not a problem.
  \item Every figure must be referred to from the text (\enquote{as shown in
        \cref{fig:learning-curves}\dots}). A figure that is not mentioned should not be
        there.
  \item \textbf{Every figure taken or adapted from another source must indicate its origin}
        and comply with its licence (see \cref{sec:quotes}). Figures without a source are
        assumed to be your own.
  \item Figures must be legible: their text should be roughly the size of the caption, and
        they should never look blurry or pixelated.
\end{itemize}

\subsection{Formats and files}\label{sec:figure-formats}

Whenever possible, use \emph{vector} formats (PDF, EPS, SVG converted to PDF) for plots and
diagrams: they keep their quality at any zoom level. Use PNG for screenshots and JPEG for
photographs, with enough resolution (at least 300~dpi at the printed size). Save your figures
in \texttt{figures/} with descriptive names, and keep the scripts that generate them in
\texttt{figures/src/}, so that they can be regenerated if the data change: the plots of this
chapter are produced by \texttt{figures/src/example\_figures.py}.

\subsection{Inserting a figure}\label{sec:inserting-figures}

Figures are inserted with \verb|\includegraphics| inside a \texttt{figure} environment, which
also contains the \verb|\caption| and the \verb|\label|. Set the size of the image relative
to the width of the text, e.g.\ \texttt{width=0.8\textbackslash textwidth}. \Cref{lst:figure} produces
\cref{fig:learning-curves}.

\begin{lstlisting}[language={[LaTeX]TeX}, caption={Inserting a figure.}, label=lst:figure]
\begin{figure}[htbp]
  \centering
  \includegraphics[width=0.85\textwidth]{example-learning-curves}
  \caption{Training and validation loss of a model. The dashed line marks the
           epoch with the lowest validation loss.}
  \label{fig:learning-curves}
\end{figure}
\end{lstlisting}

\begin{figure}[htbp]
  \centering
  \includegraphics[width=0.85\textwidth]{example-learning-curves}
  \caption{Training and validation loss of a model. The dashed line marks the epoch with
           the lowest validation loss.}
  \label{fig:learning-curves}
\end{figure}

The optional argument \texttt{[htbp]} lets \LaTeX{} place the figure \emph{here}, at the
\emph{top} or \emph{bottom} of a page, or on a separate \emph{page}, wherever it fits best.
Do not fight this behaviour: floats moving to the next page are normal, and you just refer
to them by number.

\subsection{Subfigures}\label{sec:subfigures}

Several related images can be grouped in a single figure with the \texttt{subfigure}
environment, each with its own caption and label, as in \cref{fig:subfigures}. You can
refer to the whole figure or to each part, e.g.\ \cref{fig:roc}.

\begin{figure}[htbp]
  \centering
  \begin{subfigure}[b]{0.45\textwidth}
    \centering
    \includegraphics[width=\textwidth]{example-roc-curve}
    \caption{ROC curves.}
    \label{fig:roc}
  \end{subfigure}\hfill
  \begin{subfigure}[b]{0.45\textwidth}
    \centering
    \includegraphics[width=\textwidth]{example-confusion-matrix}
    \caption{Confusion matrix of the proposed model.}
    \label{fig:confusion}
  \end{subfigure}
  \caption{Example of a figure with two subfigures.}
  \label{fig:subfigures}
\end{figure}

\subsection{Diagrams drawn with TikZ}\label{sec:tikz}

Diagrams can also be drawn directly in \LaTeX{} with the TikZ package, so that their fonts
match the text exactly. \Cref{fig:pipeline} is an example. Alternatively, draw diagrams with
tools such as draw.io (\url{https://app.diagrams.net}) or PlantUML
(\url{https://plantuml.com}) and export them as PDF.

\begin{figure}[htbp]
  \centering
  \begin{tikzpicture}[
      node distance=0.6cm,
      box/.style={draw=tfmaccent, fill=tfmaccent!8, rounded corners=2pt, align=center,
                  minimum height=1cm, minimum width=2.9cm, font=\small},
      arrow/.style={-{Stealth[length=2.5mm]}, thick, tfmaccent}]
    \node[box] (data) {Raw data};
    \node[box, right=of data] (prep) {Pre-processing};
    \node[box, right=of prep] (model) {Model training};
    \node[box, right=of model] (eval) {Evaluation};
    \draw[arrow] (data) -- (prep);
    \draw[arrow] (prep) -- (model);
    \draw[arrow] (model) -- (eval);
    \draw[arrow, dashed] (eval.south) -- ++(0,-0.7) -| node[pos=0.25, below=1pt,
      font=\footnotesize\itshape, text=tfmgray] {hyperparameter tuning} (model.south);
  \end{tikzpicture}
  \caption{A typical machine learning pipeline drawn with TikZ.}
  \label{fig:pipeline}
\end{figure}

\section{Tables}\label{sec:tables}

Tables are very common in a thesis, especially to report experimental results. The rules for
figures apply to tables too: every table has a number, a caption and a reference from the
text. In addition:

\begin{itemize}
  \item The caption of a table goes \emph{above} it (the template does it automatically
        when \verb|\caption| is written before the \texttt{tabular}).
  \item All tables share the same format, with three horizontal rules (\verb|\toprule|,
        \verb|\midrule| and \verb|\bottomrule| from the \texttt{booktabs} package), column
        headings in bold and no vertical lines.
  \item Numbers are aligned on the decimal point and reported with a sensible number of
        decimal places.
\end{itemize}

Tables are created with the \texttt{tabular} environment inside a \texttt{table} float. The
argument of \texttt{tabular} gives the alignment of each column: \texttt{l} (left),
\texttt{c} (centre), \texttt{r} (right) or \texttt{p\{5cm\}} (fixed width with line wrapping).
The \texttt{tabularx} environment adds the \texttt{X} column type, which takes all the
remaining width, very convenient for columns with long text. Columns are separated with
\verb|&| and rows end with \verb|\\|. \Cref{lst:table} produces \cref{tab:example}.

\begin{lstlisting}[language={[LaTeX]TeX}, caption={Creating a table.}, label=lst:table]
\begin{table}[htbp]
  \centering
  \caption{Datasets used in the experiments.}
  \label{tab:example}
  \begin{tabular}{@{}l r r S[table-format=2.1]@{}}
    \toprule
    \textbf{Dataset} & \textbf{Instances} & \textbf{Features} & {\textbf{Missing (\%)}} \\
    \midrule
    Hospital A & 12,480 & 54  & 3.2  \\
    Hospital B &  8,915 & 54  & 11.7 \\
    Public     & 70,000 & 112 & 0.4  \\
    \bottomrule
  \end{tabular}
\end{table}
\end{lstlisting}

\begin{table}[htbp]
  \centering
  \caption{Datasets used in the experiments.}
  \label{tab:example}
  \begin{tabular}{@{}l r r S[table-format=2.1]@{}}
    \toprule
    \textbf{Dataset} & \textbf{Instances} & \textbf{Features} & {\textbf{Missing (\%)}} \\
    \midrule
    Hospital A & 12,480 & 54  & 3.2  \\
    Hospital B &  8,915 & 54  & 11.7 \\
    Public     & 70,000 & 112 & 0.4  \\
    \bottomrule
  \end{tabular}
\end{table}

The \texttt{S} column type of the \texttt{siunitx} package aligns numbers on the decimal
point, as in the last column of \cref{tab:example}; headings in \texttt{S} columns must be
enclosed in braces. The same package typesets units consistently: \verb|\SI{8}{\giga\byte}|
gives \SI{8}{\giga\byte} and \verb|\num{0.05}| gives \num{0.05}. \Cref{tab:results} in
\cref{ch:results} shows a complete results table, and \cref{tab:hardware} shows a table
stored in its own file (\texttt{tables/hardware-requirements.tex}) and included with
\verb|\input|, which is handy for long tables or tables generated by a program.

\section{Equations}\label{sec:equations}

\LaTeX{} excels at typesetting mathematics. There are two cases:

\begin{itemize}
  \item \emph{Inline} expressions, within a paragraph, are delimited by dollar signs:
        \texttt{\$f(n) = n\textasciicircum 2 + 1\$} produces $f(n) = n^2 + 1$.
  \item \emph{Displayed} equations, on their own line, numbered and labelled so that they
        can be referenced. They use the \texttt{equation} environment, or \texttt{align}
        for several aligned equations.
\end{itemize}

\Cref{lst:equation} produces \cref{eq:cross-entropy}, the binary cross-entropy loss.

\begin{lstlisting}[language={[LaTeX]TeX}, caption={A displayed equation.}, label=lst:equation]
\begin{equation}
  \mathcal{L}(\bm{\theta}) = -\frac{1}{N} \sum_{i=1}^{N} \left[ y_i \log \hat{y}_i
                             + (1 - y_i) \log (1 - \hat{y}_i) \right]
  \label{eq:cross-entropy}
\end{equation}
\end{lstlisting}

\begin{equation}
  \mathcal{L}(\bm{\theta}) = -\frac{1}{N} \sum_{i=1}^{N} \left[ y_i \log \hat{y}_i
                             + (1 - y_i) \log (1 - \hat{y}_i) \right]
  \label{eq:cross-entropy}
\end{equation}

Equations are part of the text, so they are punctuated like any other sentence, and every
symbol must be defined the first time it appears: here $N$ is the number of examples, $y_i$
the true label of the $i$-th example and $\hat{y}_i$ the predicted probability. The
Wikibooks guide (\url{https://en.wikibooks.org/wiki/LaTeX/Mathematics}) answers most
questions about mathematical notation.

\section{Algorithms}\label{sec:algorithms}

Algorithms that are used or proposed should be described in pseudocode, with the same syntax
and format throughout the thesis. The \texttt{algorithm2e} package provides the
\texttt{algorithm} environment, with its own caption and label, and commands for
conditionals, loops, etc. \Cref{lst:algorithm} produces \cref{alg:gradient-descent}.

\begin{algorithm}[htbp]
  \caption{Mini-batch gradient descent.}
  \label{alg:gradient-descent}
  \KwIn{Training set $\mathcal{D}$, learning rate $\eta$, number of epochs $E$}
  \KwOut{Trained parameters $\bm{\theta}$}
  Initialise $\bm{\theta}$ randomly\;
  \For{$e \leftarrow 1$ \KwTo $E$}{
    Shuffle $\mathcal{D}$\;
    \ForEach{mini-batch $\mathcal{B} \subset \mathcal{D}$}{
      $\bm{\theta} \leftarrow \bm{\theta} - \eta \, \nabla_{\bm{\theta}} \mathcal{L}(\bm{\theta}; \mathcal{B})$\;
    }
    \If{the validation loss has not improved for 10 epochs}{
      \textbf{break}\;
    }
  }
  \Return{$\bm{\theta}$}\;
\end{algorithm}

\begin{lstlisting}[language={[LaTeX]TeX}, caption={Describing an algorithm.}, label=lst:algorithm]
\begin{algorithm}[htbp]
  \caption{Mini-batch gradient descent.}
  \label{alg:gradient-descent}
  \KwIn{Training set $\mathcal{D}$, learning rate $\eta$, number of epochs $E$}
  \KwOut{Trained parameters $\bm{\theta}$}
  Initialise $\bm{\theta}$ randomly\;
  \For{$e \leftarrow 1$ \KwTo $E$}{
    Shuffle $\mathcal{D}$\;
    \ForEach{mini-batch $\mathcal{B} \subset \mathcal{D}$}{
      $\bm{\theta} \leftarrow \bm{\theta} - \eta \, \nabla_{\bm{\theta}} \mathcal{L}(\bm{\theta}; \mathcal{B})$\;
    }
    \If{the validation loss has not improved for 10 epochs}{
      \textbf{break}\;
    }
  }
  \Return{$\bm{\theta}$}\;
\end{algorithm}
\end{lstlisting}

Note that each instruction ends with \verb|\;| and that the caption of an algorithm appears
at the top. The documentation of \texttt{algorithm2e}
(\url{https://ctan.org/pkg/algorithm2e}) describes all the available commands, although a
small subset is usually enough.

\section{Source code}\label{sec:code}

Besides pseudocode, a thesis may include listings of real code. All the listings in this
chapter have been produced with the \texttt{lstlisting} environment of the
\texttt{listings} package, which takes three main parameters: \texttt{language}, which
enables syntax highlighting; \texttt{caption}; and \texttt{label}. \Cref{lst:python} shows a
Python example. Complete source files can be included directly from disk, without copying
their content, with the \verb|\lstinputlisting| command:

\begin{lstlisting}[language={[LaTeX]TeX}, numbers=none]
\lstinputlisting[language=Python, caption={...}, label=lst:...]{code/file.py}
\end{lstlisting}

\begin{lstlisting}[language=Python, caption={Training a classifier with scikit-learn.}, label=lst:python]
from sklearn.ensemble import RandomForestClassifier
from sklearn.model_selection import cross_val_score

# 10-fold cross-validation of a random forest
model = RandomForestClassifier(n_estimators=500, random_state=2026)
scores = cross_val_score(model, X, y, cv=10, scoring="f1")
print(f"F1 = {scores.mean():.3f} ± {scores.std():.3f}")
\end{lstlisting}

The list of supported languages, and how to define new ones, can be found at
\url{https://en.wikibooks.org/wiki/LaTeX/Source_Code_Listings}. Include only the code that
helps to understand the thesis; the complete code belongs to a repository referenced from
the text (see \cref{sec:reproducibility}).

\section{Acronyms}\label{sec:acronyms}

Acronyms are defined once in \texttt{frontmatter/acronyms.tex}, as in
\cref{lst:acronym}, and used in the text with \verb|\gls{key}|.

\begin{lstlisting}[language={[LaTeX]TeX}, caption={Defining an acronym.}, label=lst:acronym]
\newacronym{gpu}{GPU}{graphics processing unit}
\end{lstlisting}

\noindent The first time an acronym is used, its long form is written followed by the short form in
brackets; afterwards only the short form appears, and every acronym used is added to the
List of Acronyms automatically. For example, the first \verb|\gls{ml}| of this chapter
produces \enquote{\gls{ml}}, and the next one produces \enquote{\gls{ml}}. The plural is
obtained with \verb|\glspl{key}|, as in \enquote{\glspl{gpu}}.

\section{Citations and references}\label{sec:citations}

During the thesis you will study books, scientific articles, theses, standards and web pages,
and you may use definitions, data or figures taken from them. In all these cases it is
essential to cite the source properly. Poor citation practice is one of the most common
reasons for a lower grade, and plagiarism, even unintentional, can lead to failing the
thesis.

\subsection{Bibliographic references}\label{sec:bibliography}

References are cited with \verb|\cite{key}|, where \texttt{key} is the identifier
of the corresponding entry in \texttt{references.bib}. The process is:

\begin{enumerate}
  \item Get the bibliographic data in BibTeX format: from Google Scholar (the
        \enquote{Cite} button, then \enquote{BibTeX}), from the publisher's page, from
        \url{https://doi.org} (which returns BibTeX for any \gls{doi}), or from a reference
        manager such as Zotero with the Better BibTeX extension, which keeps
        \texttt{references.bib} up to date automatically.
  \item Add the entry to \texttt{references.bib} and give it an identifier that is easy to
        remember, e.g.\ \texttt{surname2024keyword}.
  \item Cite it in the text with \verb|\cite{identifier}|.
\end{enumerate}

Each type of work has its own type of entry: \texttt{@article}, \texttt{@book},
\texttt{@inproceedings}, \texttt{@thesis}, \texttt{@online}, \texttt{@software},
\texttt{@dataset}\dots{} In general, every entry needs at least \texttt{author},
\texttt{title} and \texttt{year} (or \texttt{date}); books also need the
\texttt{publisher}; articles, the \texttt{journal}, \texttt{volume}, \texttt{number} and
\texttt{pages}; web resources, the \texttt{url} and the date when you accessed them
(\texttt{urldate}). Always add the \texttt{doi} when the work has one. Automatically
generated entries are often incomplete or contain errors, so review them carefully.
\Cref{lst:bibtex} shows four entries of different types, which you can also find in
\texttt{references.bib}.

\begin{lstlisting}[language={[LaTeX]TeX}, caption={BibLaTeX entries of different types.}, label=lst:bibtex]
@article{charte2020review,
  title   = {A Comprehensive and Didactic Review on Multilabel
             Learning Software Tools},
  author  = {Charte, Francisco},
  journal = {IEEE Access},
  volume  = {8},
  pages   = {50330--50354},
  year    = {2020},
  doi     = {10.1109/ACCESS.2020.2979787}
}
@book{herrera2016multilabel,
  title     = {Multilabel Classification: Problem Analysis, Metrics
               and Techniques},
  author    = {Herrera, Francisco and Charte, Francisco and
               Rivera, Antonio J. and del Jesus, Mar{\'\i}a J.},
  year      = {2016},
  publisher = {Springer},
  isbn      = {978-3-319-41110-1},
  doi       = {10.1007/978-3-319-41111-8}
}
@inproceedings{charte2013imbalance,
  title     = {A First Approach to Deal with Imbalance in
               Multi-label Datasets},
  author    = {Charte, Francisco and Rivera, Antonio and
               del Jesus, Mar{\'\i}a Jos{\'e} and Herrera, Francisco},
  booktitle = {Hybrid Artificial Intelligent Systems (HAIS 2013)},
  pages     = {150--160},
  year      = {2013},
  doi       = {10.1007/978-3-642-40846-5_16}
}
@online{charte2019distribution,
  author  = {Charte, Francisco},
  title   = {Cómo analizar la distribución de los datos con {R}},
  year    = {2019},
  url     = {https://fcharte.com/tutoriales/20170114-DistribucionDatosR/},
  urldate = {2020-09-30},
  langid  = {spanish}
}
\end{lstlisting}

With the entries ready, the paragraph in \cref{lst:citations} produces the text that follows
it, and the complete references are added automatically, correctly formatted, to the
References section at the end of the thesis. Only cited works appear there.

\begin{lstlisting}[language={[LaTeX]TeX}, caption={A paragraph citing several works.}, label=lst:citations]
As explained in~\cite{charte2020review} and extended in the
book~\cite{herrera2016multilabel}, the techniques presented
in~\cite{charte2013imbalance} can be put into practice by following
the tutorial in~\cite{charte2019distribution}.
\end{lstlisting}

\begin{quotation}
  As explained in~\cite{charte2020review} and extended in the
  book~\cite{herrera2016multilabel}, the techniques presented in~\cite{charte2013imbalance}
  can be put into practice by following the tutorial in~\cite{charte2019distribution}.
\end{quotation}

The citation style is chosen with the \texttt{bibstyle} class option (\cref{tab:options}).
Numeric styles are common in engineering and computer science, while author--year styles
such as APA are common in health and social sciences. Ask your supervisor which one is
expected in your field. With author--year styles, use \verb|\textcite{key}| to cite the
authors as part of the sentence, e.g.\ \enquote{\textcite{pedregosa2011scikit} showed
that\dots}, and \verb|\parencite{key}| for citations in parentheses.

\subsection{Quotations and third-party material}\label{sec:quotes}

If, besides using a work as a source, you reproduce part of it literally, follow these rules:

\begin{itemize}
  \item Only short fragments of a work can be quoted literally, usually no more than a
        paragraph and never whole pages.
  \item Short quotations go in the text between quotation marks, with \verb|\enquote{...}|;
        longer quotations are set apart with the \texttt{displayquote} environment:
        \begin{displayquote}
          This is how a long literal quotation looks: it is indented and separated from the
          rest of the text, so that the reader can clearly tell the words of the author
          apart from yours.~\cite{charte2020review}
        \end{displayquote}
  \item Every quotation must be accompanied by its bibliographic citation, so that its
        origin can be identified unambiguously.
  \item Do not include illustrations, diagrams or figures whose rights of use are not
        clear. Material under a Creative Commons licence can be included by indicating its
        source and licence; otherwise you need the owner's permission. The alternative is
        to draw your own figure, indicating that it is based on an existing source
        (\enquote{adapted from~\cite{herrera2016multilabel}}).
\end{itemize}

Remember that your thesis will probably be published in your university's institutional
repository, where anyone can read it, and that it will be checked with plagiarism detection
tools. Be honest, and never use the work of others without citing it.

\section{Special characters}\label{sec:special-characters}

Some characters have a special meaning in \LaTeX: \verb|\| starts commands, \verb|$| delimits
mathematics, \verb|%| starts a comment, \verb|&| separates table columns, \verb|#| denotes
parameters, \verb|_| and \verb|^| create subscripts and superscripts, and \verb|{ }| group
content. To print them, precede them with a backslash: \verb|\%|, \verb|\$|, \verb|\&|,
\verb|\#|, \verb|\_|, \verb|\{| and \verb|\}| produce \%, \$, \&, \#, \_, \{ and \}. The
backslash itself is written with \verb|\textbackslash| (\textbackslash), because
\verb|\\| means a line break. Accented letters such as á, é, ñ or ç can be typed directly.
If you have trouble with any character, see
\url{https://en.wikibooks.org/wiki/LaTeX/Special_Characters}.

\section{Working with your supervisors}\label{sec:review}

When you send a draft to your supervisors, compile it with the \texttt{review} option: every
line is numbered, so comments such as \enquote{page 12, line 34} are unambiguous, and the
watermark prevents a draft from being mistaken for the final version. While writing, you can
leave notes for yourself or your supervisors with \verb|\todo{text}|, which appears in the
margin\todo{This is a to-do note.}, or with \verb|\todo[inline]{text}|:

\todo[inline]{This is an inline to-do note. All notes disappear with the \texttt{final}
option.}

Use a version control system such as Git from the very first day (Overleaf can synchronise
with GitHub), so that you never lose work and can always go back to a previous version.

\section{Structure of a Master's thesis}\label{sec:tfm-structure}

Once your data are set, the next step is writing the content. The structure depends on the
type of thesis and on the regulations of your programme, so take the following as a
recommendation that can be adapted to each case. Each chapter is stored in its own
\texttt{.tex} file and included from \texttt{main.tex} with \verb|\input|: add, remove or
rename the files as you need.

\subsection{Research thesis}\label{sec:research-thesis}

This is the most common type of Master's thesis, and the one this template is prepared for.
Its chapters follow the usual structure of a scientific article:

\begin{enumerate}
  \item \textbf{Introduction} --- \texttt{01-introduction.tex}. A high-level overview of
        the work: motivation, importance of the problem, how it has been addressed so far
        and what this thesis contributes, without going into detail. It usually ends with
        the structure of the document.
  \item \textbf{Background and Related Work} --- \texttt{02-background.tex}. Shows that you
        have studied the topic: the problem in detail, the techniques used so far and those
        you propose. The bibliography is fundamental in this chapter.
  \item \textbf{Objectives} --- \texttt{03-objectives.tex}. The general objective, broken
        down into specific ones, and, if appropriate, the hypothesis or research questions.
        Just by reading this chapter, it should be clear what has been done.
  \item \textbf{Materials and Methods} --- \texttt{04-materials-and-methods.tex}. The data
        and resources used and the techniques applied, in enough detail to reproduce the
        work, including the experimental setup, the evaluation and, when relevant, the
        ethical and legal aspects.
  \item \textbf{Results} --- \texttt{05-results.tex}. All the results obtained, usually in
        tables and figures, compared with other methods whenever possible.
  \item \textbf{Discussion} --- \texttt{06-discussion.tex}. The interpretation of the
        results, their comparison with previous work, their limitations and their
        implications. Some fields merge it with the previous chapter.
  \item \textbf{Conclusions and Future Work} --- \texttt{07-conclusions.tex}. A summary of
        the work and results, an assessment of whether the objectives of the approved
        proposal have been achieved, and the lines of future work.
\end{enumerate}

\subsection{Professional or engineering project}\label{sec:project-thesis}

In professionally oriented Master's degrees, the thesis is often a project: the development
of a software system, a product or a solution for an organisation. Its usual chapters are:

\begin{enumerate}
  \item \textbf{Introduction}, as described above.
  \item \textbf{Planning and Budget}: time and cost planning of the project and the
        resources it needs; the methodology used to manage it (e.g.\ Scrum); the
        estimation of the size and effort of the software, from which the schedule (usually
        a Gantt chart, see \cref{fig:gantt}) and the budget are derived; and the legal and
        technical regulations followed.
  \item \textbf{Requirements Analysis}: the interviews with the (real or hypothetical)
        users or clients, from which the functional and non-functional requirements are
        identified, analysed, specified (use case diagrams, user stories, domain model\dots)
        and validated, following the practices recommended in the ISO/IEC/IEEE 29148
        standard.
  \item \textbf{System Design}: the detailed design of the system core (architecture,
        package, class and sequence diagrams), its user interface (wireframes or mockups)
        and its data (entity--relationship diagrams or equivalent), including the
        architectural and design patterns used.
  \item \textbf{Implementation}: the evaluation of the alternative languages, libraries and
        tools and the justified choice among them; the implementation iterations; the
        structure of the project; the coding and documentation style guides; and the build
        and configuration process.
  \item \textbf{Testing}: the tools used and the unit, integration, regression and
        acceptance tests carried out.
  \item \textbf{Deployment and Maintenance}: how the system is deployed and put into
        operation, its dependencies on external hosting or services, and the operations
        needed to maintain it.
  \item \textbf{Conclusions and Future Work}, as described above.
  \item \textbf{Installation and User Manuals}, usually as appendices, written for the final
        users rather than for developers.
\end{enumerate}

Projects of this type need diagrams of several kinds. Tools such as PlantUML, which define
diagrams as text files that can be kept under version control, are highly recommended.

\section{Writing advice}\label{sec:writing-advice}

Although it has nothing to do with \LaTeX{}, the quality of the writing deserves a final
word. For the examination board and for any future reader, the quality of the text is a
reflection of the quality of the work. A very good piece of work can be penalised because of
poor writing, while a less brilliant one with an impeccable presentation may be valued more
highly. Without trying to be exhaustive, keep in mind the following:

\begin{itemize}
  \item \textbf{Spelling:} choose British or American English (option \texttt{english})
        and be consistent: \textit{analyse}/\textit{analyze},
        \textit{modelling}/\textit{modeling}, \textit{colour}/\textit{color}. Use a spell
        checker, such as the one integrated in Overleaf, set to the same variety.
  \item \textbf{Grammar and tense:} use the present tense for established knowledge and for
        what the thesis contains (\enquote{\cref{ch:results} presents\dots}) and the past
        tense for what you did (\enquote{the data were collected\dots}). Avoid
        contractions (\textit{don't}, \textit{it's}) in formal writing.
  \item \textbf{Common pitfalls for Spanish speakers:} watch out for false friends:
        \textit{actually} means \enquote{en realidad} (not \enquote{actualmente}, which is
        \textit{currently}), \textit{eventually} means \enquote{finalmente}, and
        \textit{sensible} means \enquote{sensato}. Do not use \enquote{the} with general
        plural nouns (\enquote{\textit{Patients} with diabetes\dots}, not \enquote{\textit{The}
        patients\dots}), adjectives never take a plural \textit{-s}, and experiments are
        \textit{conducted} or \textit{carried out}, not \textit{made}. Acronyms do take an
        \textit{s} in the plural in English (\textit{GPUs}).
  \item \textbf{Vocabulary:} use precise, specialised terms, which show that you know the
        field, and vary your vocabulary to avoid repeating the same words. A thesis is a
        formal document: write \enquote{as described in\dots} rather than \enquote{as we
        said before\dots}.
  \item \textbf{Structure:} the thesis as a whole, each chapter and each section must read
        as a sequence of connected ideas. Nothing is more confusing than a text in which
        each paragraph jumps to a different topic. Plan the content of each chapter and
        section before writing it so that, as in good software design, cohesion is maximised
        and coupling minimised.
  \item \textbf{Generative AI:} if you use AI tools to improve your writing, follow the
        rules of your university and declare their use (see
        \texttt{frontmatter/ai-statement.tex}). You remain the author and the person
        responsible for every sentence of the thesis.
\end{itemize}

One last piece of advice, completely outside the scope of these instructions: the best way to
learn to write well is to read a lot, especially good scientific articles in your field and
good literature.

\section{Getting this template}\label{sec:getting-template}

This template is publicly available at
\url{https://github.com/thisisjosepablo/TFMLatexTemplate}, where you can download the latest
release, open it directly in Overleaf, report problems or contribute improvements, such as
settings that reproduce the style of your own university. It is released into the public
domain under the CC0~1.0 licence, so you can use it however you like, without any condition,
but also without any warranty.
